% GENERATED FILE. Edit canonical lessons, then run npm run generate:books.
% canonical-source-hash: fnv1a64:371633eeb996768c
% canonical-lessons: GU-C194-naku, GU-C194-signal, GU-C194-same, GU-C194-baju, GU-C194-olangvu, GU-R194-review-189-194

\chapter{Crossing at the corner}
\label{ch:gu-crossing-at-the-corner}
\hlchaptermodality{194}

\begin{chapteropening}
\textbf{By the end of this chapter:} \emph{I can say where a place stands: at the corner, opposite or beside.}

Complete the last lesson of chapter 194: I can say where a place stands: at the corner, opposite or beside.
\end{chapteropening}

\section[\texorpdfstring{\gu{નાકું} (nākũ)}{નાકું (nākũ)}]{\gu{નાકું} (nākũ) — street corner, end of a road}

\label{lesson:GU-C194-naku}



\begin{quote}
Before the new one: say the Gujarati for straight, then the Gujarati for to turn.
\end{quote}

\subsection*{You'll want to know: \gu{નાકું}}
\textbf{\gu{નાકું}} — \emph{nākũ} — \textquotedblleft{}street corner, end of a road\textquotedblright{}.

\begin{quote}
\textbf{\gu{દુકાન} \gu{ગલીના} \gu{નાકે} \gu{છે}.} — \emph{dukān galīnā nāke che} — \textquotedblleft{}The shop is at the end of the lane.\textquotedblright{}
\end{quote}

A neuter noun; \textbf{\gu{નાકે}}, \emph{nāke}, 'at the corner'.

\begin{cousinweb}
\textbf{\gu{નાકું}} (\emph{nākũ}) is a tadbhava. CDIAL lists it, glossed \textquotedblleft{}eye of a needle, end of a road, passage\textquotedblright{}, under \emph{*nākka}, \textquotedblleft{}like a nose, projecting\textquotedblright{}, a reconstructed form (CDIAL 7037).
\end{cousinweb}

\begin{grammarlens}[title={What you've built}]
One more word for everyday life at A2.
\end{grammarlens}

\subsection*{Your turn}
\begin{itemize}
\raggedright
  \item \emph{Say it:} \emph{nākũ}
  \item \emph{Say it:} \emph{nākũ}, once more
  \item \emph{Say it:} \emph{dukān galīnā nāke che} — The shop is at the end of the lane.
  \item \emph{Recall:} say the Gujarati for straight, then the Gujarati for to turn, then say \emph{nākũ} again
\end{itemize}

\begin{tcolorbox}[breakable,colback=teal!4,colframe=teal!35!black,title={Before you move on}]
What does \gu{નાકું} mean? (\textquotedblleft{}Street corner, end of a road\textquotedblright{}.) Say it once more.
\end{tcolorbox}
\section[\texorpdfstring{\gu{સિગ્નલ} (signal)}{સિગ્નલ (signal)}]{\gu{સિગ્નલ} (signal) — traffic signal}

\label{lesson:GU-C194-signal}



\begin{quote}
Before the new one: say the Gujarati for to turn, then the Gujarati for street corner.
\end{quote}

\subsection*{You'll want to know: \gu{સિગ્નલ}}
\textbf{\gu{સિગ્નલ}} — \emph{signal} — \textquotedblleft{}traffic signal\textquotedblright{}.

\begin{quote}
\textbf{\gu{સિગ્નલ} \gu{લાલ} \gu{છે}.} — \emph{signal lāl che} — \textquotedblleft{}The signal is red.\textquotedblright{}
\end{quote}

\begin{cousinweb}
\textbf{\gu{સિગ્નલ}} (\emph{signal}) is a loan from English \emph{signal}.
\end{cousinweb}

\begin{grammarlens}[title={What you've built}]
One more word for everyday life at A2.
\end{grammarlens}

\subsection*{Your turn}
\begin{itemize}
\raggedright
  \item \emph{Say it:} \emph{signal}
  \item \emph{Say it:} \emph{signal}, once more
  \item \emph{Say it:} \emph{signal lāl che} — The signal is red.
  \item \emph{Recall:} say the Gujarati for to turn, then the Gujarati for street corner, then say \emph{signal} again
\end{itemize}

\begin{tcolorbox}[breakable,colback=teal!4,colframe=teal!35!black,title={Before you move on}]
What does \gu{સિગ્નલ} mean? (\textquotedblleft{}Traffic signal\textquotedblright{}.) Say it once more.
\end{tcolorbox}
\section[\texorpdfstring{\gu{સામે} (sāme)}{સામે (sāme)}]{\gu{સામે} (sāme) — opposite, facing}

\label{lesson:GU-C194-same}



\begin{quote}
Before the new one: say the Gujarati for street corner, then the Gujarati for traffic signal.
\end{quote}

\subsection*{You'll want to know: \gu{સામે}}
\textbf{\gu{સામે}} — \emph{sāme} — \textquotedblleft{}opposite, facing\textquotedblright{}.

\begin{quote}
\textbf{\gu{બેંક} \gu{સ્ટેશનની} \gu{સામે} \gu{છે}.} — \emph{bẽk sṭeśannī sāme che} — \textquotedblleft{}The bank is opposite the station.\textquotedblright{}
\end{quote}

\begin{cousinweb}
\textbf{\gu{સામે}} (\emph{sāme}) is a tadbhava. It comes from Sanskrit \emph{saṁmukhá}, \textquotedblleft{}facing, present\textquotedblright{}, and CDIAL lists Gujarati \emph{sāmũ}, \textquotedblleft{}facing\textquotedblright{} (CDIAL 12982).
\end{cousinweb}

\begin{grammarlens}[title={What you've built}]
One more word for everyday life at A2.
\end{grammarlens}

\subsection*{Your turn}
\begin{itemize}
\raggedright
  \item \emph{Say it:} \emph{sāme}
  \item \emph{Say it:} \emph{sāme}, once more
  \item \emph{Say it:} \emph{bẽk sṭeśannī sāme che} — The bank is opposite the station.
  \item \emph{Recall:} say the Gujarati for street corner, then the Gujarati for traffic signal, then say \emph{sāme} again
\end{itemize}

\begin{tcolorbox}[breakable,colback=teal!4,colframe=teal!35!black,title={Before you move on}]
What does \gu{સામે} mean? (\textquotedblleft{}Opposite, facing\textquotedblright{}.) Say it once more.
\end{tcolorbox}
\section[\texorpdfstring{\gu{બાજુ} (bāju)}{બાજુ (bāju)}]{\gu{બાજુ} (bāju) — side}

\label{lesson:GU-C194-baju}



\begin{quote}
Before the new one: say the Gujarati for traffic signal, then the Gujarati for opposite.
\end{quote}

\subsection*{You'll want to know: \gu{બાજુ}}
\textbf{\gu{બાજુ}} — \emph{bāju} — \textquotedblleft{}side\textquotedblright{}.

\begin{quote}
\textbf{\gu{હોટેલ} \gu{મંદિરની} \gu{બાજુમાં} \gu{છે}.} — \emph{hoṭel mandirnī bājumā̃ che} — \textquotedblleft{}The hotel is next to the temple.\textquotedblright{}
\end{quote}

A feminine noun.

\begin{cousinweb}
\textbf{\gu{બાજુ}} (\emph{bāju}) is a loan, usually traced to Persian \emph{bāzū}, \textquotedblleft{}arm\textquotedblright{}, and so \textquotedblleft{}side\textquotedblright{}. CDIAL has no entry for it.
\end{cousinweb}

\begin{grammarlens}[title={What you've built}]
One more word for everyday life at A2.
\end{grammarlens}

\subsection*{Your turn}
\begin{itemize}
\raggedright
  \item \emph{Say it:} \emph{bāju}
  \item \emph{Say it:} \emph{bāju}, once more
  \item \emph{Say it:} \emph{hoṭel mandirnī bājumā̃ che} — The hotel is next to the temple.
  \item \emph{Recall:} say the Gujarati for traffic signal, then the Gujarati for opposite, then say \emph{bāju} again
\end{itemize}

\begin{tcolorbox}[breakable,colback=teal!4,colframe=teal!35!black,title={Before you move on}]
What does \gu{બાજુ} mean? (\textquotedblleft{}Side\textquotedblright{}.) Say it once more.
\end{tcolorbox}
\section[\texorpdfstring{\gu{ઓળંગવું} (oḷangvũ)}{ઓળંગવું (oḷangvũ)}]{\gu{ઓળંગવું} (oḷangvũ) — to cross}

\label{lesson:GU-C194-olangvu}



\begin{quote}
Before the new one: say the Gujarati for opposite, then the Gujarati for side.
\end{quote}

\subsection*{You'll want to know: \gu{ઓળંગવું}}
\textbf{\gu{ઓળંગવું}} — \emph{oḷangvũ} — \textquotedblleft{}to cross\textquotedblright{}.

\begin{quote}
\textbf{\gu{બાળકો} \gu{સિગ્નલ} \gu{પાસે} \gu{રસ્તો} \gu{ઓળંગે} \gu{છે}.} — \emph{bāḷko signal pāse rasto oḷange che} — \textquotedblleft{}Children cross the road at the signal.\textquotedblright{}
\end{quote}

\begin{cousinweb}
\textbf{\gu{ઓળંગવું}} (\emph{oḷangvũ}) is a tadbhava. CDIAL lists it as \emph{oḷãghvũ}, \textquotedblleft{}to cross, violate\textquotedblright{}, under \emph{*avalaṅghati}, a reconstructed form built on \emph{laṅgh}, \textquotedblleft{}to leap\textquotedblright{} (CDIAL 824).
\end{cousinweb}

\begin{grammarlens}[title={What you've built}]
One more word for everyday life at A2.
\end{grammarlens}

\subsection*{Your turn}
\begin{itemize}
\raggedright
  \item \emph{Say it:} \emph{oḷangvũ}
  \item \emph{Say it:} \emph{oḷangvũ}, once more
  \item \emph{Say it:} \emph{bāḷko signal pāse rasto oḷange che} — Children cross the road at the signal.
  \item \emph{Recall:} say the Gujarati for opposite, then the Gujarati for side, then say \emph{oḷangvũ} again
\end{itemize}

\begin{tcolorbox}[breakable,colback=teal!4,colframe=teal!35!black,title={Before you move on}]
What does \gu{ઓળંગવું} mean? (\textquotedblleft{}To cross\textquotedblright{}.) Say it once more.
\end{tcolorbox}
\section[(dialogue)]{Words from chapters 189–194 — a review}

\label{lesson:GU-R194-review-189-194}



\begin{quote}
Say the Gujarati for to wrap and to order.
\end{quote}

\subsection*{Your turn}
Say these aloud:

\begin{itemize}
\raggedright
  \item to wrap, to order, to run out, to be left over, to accept
  \item bank account, to deposit, cheque, to cash a cheque, to borrow
  \item to lend, interest on money, savings, profit, loss
  \item reckoning, total, amount, bill, to waste
  \item direction, left, right-hand, straight, to turn
  \item street corner, traffic signal, opposite, side, to cross
\end{itemize}

\begin{tcolorbox}[breakable,colback=teal!4,colframe=teal!35!black,title={Before you move on}]
To wrap? (\textbf{\gu{લપેટવું}}.) To order? (\textbf{\gu{મંગાવવું}}.) To run out? (\textbf{\gu{ખૂટવું}}.) To be left over? (\textbf{\gu{બચવું}}.) To accept? (\textbf{\gu{સ્વીકારવું}}.) Bank account? (\textbf{\gu{ખાતું}}.) To deposit? (\textbf{\gu{જમા} \gu{કરવું}}.) Cheque? (\textbf{\gu{ચેક}}.) To cash a cheque? (\textbf{\gu{વટાવવું}}.) To borrow? (\textbf{\gu{ઉધાર} \gu{લેવું}}.) To lend? (\textbf{\gu{ધીરવું}}.) Interest on money? (\textbf{\gu{વ્યાજ}}.) Savings? (\textbf{\gu{બચત}}.) Profit? (\textbf{\gu{નફો}}.) Loss? (\textbf{\gu{ખોટ}}.) Reckoning? (\textbf{\gu{હિસાબ}}.) Total? (\textbf{\gu{કુલ}}.) Amount? (\textbf{\gu{રકમ}}.) Bill? (\textbf{\gu{બિલ}}.) To waste? (\textbf{\gu{વેડફવું}}.) Direction? (\textbf{\gu{દિશા}}.) Left? (\textbf{\gu{ડાબું}}.) Right-hand? (\textbf{\gu{જમણું}}.) Straight? (\textbf{\gu{સીધું}}.) To turn? (\textbf{\gu{વળવું}}.) Street corner? (\textbf{\gu{નાકું}}.) Traffic signal? (\textbf{\gu{સિગ્નલ}}.) Opposite? (\textbf{\gu{સામે}}.) Side? (\textbf{\gu{બાજુ}}.) To cross? (\textbf{\gu{ઓળંગવું}}.)
\end{tcolorbox}
